\documentclass[letterpaper]{article} % DO NOT CHANGE THIS
\usepackage[]{aaai2027}  % DO NOT CHANGE THIS
% The serif, sans-serif, and monospaced fonts are loaded automatically by
% aaai2027.sty (newtxtext, helvet, courier). DO NOT add \usepackage{times},
% \usepackage{helvet}, \usepackage{courier}, or any other font package.
\usepackage[hyphens]{url}  % DO NOT CHANGE THIS
\usepackage{graphicx} % DO NOT CHANGE THIS
\urlstyle{rm} % DO NOT CHANGE THIS
\def\UrlFont{\rm}  % DO NOT CHANGE THIS
\usepackage{natbib}  % DO NOT CHANGE THIS AND DO NOT ADD ANY OPTIONS TO IT
\usepackage{caption} % DO NOT CHANGE THIS AND DO NOT ADD ANY OPTIONS TO IT
\frenchspacing  % DO NOT CHANGE THIS
%
% These are recommended to typeset algorithms but not required. See the subsubsection on algorithms. Remove them if you don't have algorithms in your paper.
\usepackage{algorithm}
\usepackage{amssymb}
\usepackage{algorithmic}
\usepackage{multirow}      % \multirow{2}{*}{...}
\usepackage{array}         
\usepackage{colortbl}      
\definecolor{ourrow}{RGB}{230, 240, 255}

%
% These are recommended to typeset listings but not required. See the subsubsection on listing. Remove this block if you don't have listings in your paper.
\usepackage{newfloat}
\usepackage{listings}
\DeclareCaptionStyle{ruled}{labelfont=normalfont,labelsep=colon,strut=off} % DO NOT CHANGE THIS
% \lstset{%
% 	basicstyle={\footnotesize\ttfamily},% footnotesize acceptable for monospace
% 	numbers=left,numberstyle=\footnotesize,xleftmargin=2em,% show line numbers, remove this entire line if you don't want the numbers.
% 	aboveskip=0pt,belowskip=0pt,%
% 	showstringspaces=false,tabsize=2,breaklines=true}
% \floatstyle{ruled}
% \newfloat{listing}{tb}{lst}{}
% \floatname{listing}{Listing}

%
% Recommended for better-looking tables
\usepackage{booktabs}
%
% Keep the \pdfinfo as shown here. There's no need
% for you to add the /Title and /Author tags.
\pdfinfo{
/TemplateVersion (2027.1)
}
\definecolor{rowgray}{HTML}{EDEDED}
\definecolor{myblue}{RGB}{155, 0, 0} 
\definecolor{rowpink}{RGB}{255, 230, 235}  % 微粉色
\definecolor{rowours}{HTML}{E8F1FA}
\usepackage{xcolor}
\definecolor{myorange}{RGB}{230, 126, 34}
\definecolor{deepred}{RGB}{139,0,0}
\usepackage{amsmath}
% DISALLOWED PACKAGES
% \usepackage{authblk} -- This package is specifically forbidden
% \usepackage{balance} -- This package is specifically forbidden
% \usepackage{CJK} -- This package is specifically forbidden
% \usepackage{float} -- This package is specifically forbidden
% \usepackage{flushend} -- This package is specifically forbidden
% \usepackage{fullpage} -- This package is specifically forbidden
% \usepackage{geometry} -- This package is specifically forbidden
% \usepackage{hyperref} -- This package is specifically forbidden
% \usepackage{navigator} -- This package is specifically forbidden
% (or any other package that embeds links such as navigator or hyperref)
% \usepackage{indentfirst} -- This package is specifically forbidden
% \usepackage{layout} -- This package is specifically forbidden
% \usepackage{multicol} -- This package is specifically forbidden
% \usepackage{nameref} -- This package is specifically forbidden
% \usepackage{savetrees} -- This package is specifically forbidden
% \usepackage{setspace} -- This package is specifically forbidden
% \usepackage{stfloats} -- This package is specifically forbidden
% \usepackage{tabu} -- This package is specifically forbidden
% \usepackage{titlesec} -- This package is specifically forbidden
% \usepackage{tocbibind} -- This package is specifically forbidden
% \usepackage{ulem} -- This package is specifically forbidden
% \usepackage{wrapfig} -- This package is specifically forbidden
% DISALLOWED COMMANDS
% \nocopyright -- Your paper will not be published if you use this command
% \addtolength -- This command may not be used
% \balance -- This command may not be used
% \baselinestretch -- Your paper will not be published if you use this command
% \clearpage -- No page breaks of any kind may be used for the final version of your paper
% \columnsep -- This command may not be used
% \newpage -- No page breaks of any kind may be used for the final version of your paper
% \pagebreak -- No page breaks of any kind may be used for the final version of your paper
% \pagestyle -- This command may not be used
% \tiny -- This is not an acceptable font size.
% \vspace{- -- No negative value may be used in proximity of a caption, figure, table, section, subsection, subsubsection, or reference
% \vskip{- -- No negative value may be used to alter spacing above or below a caption, figure, table, section, subsection, subsubsection, or reference

\setcounter{secnumdepth}{2} %May be changed to 1 or 2 if section numbers are desired.

% The file aaai2027.sty is the style file for AAAI Press
% proceedings, working notes, and technical reports.
%

% Title

% Your title must be in mixed case, not sentence case.
% That means all verbs (including short verbs like be, is, using,and go),
% nouns, adverbs, adjectives should be capitalized, including both words in hyphenated terms, while
% articles, conjunctions, and prepositions are lower case unless they
% directly follow a colon or long dash
\title{\raisebox{-0.3\height}{\includegraphics[height=1.3cm]{images/g-carl-logo-v5-transparent.png}} G-CARL: Grounded Checklist-Aligned Reward Learning for Patient-Oriented Medical Report Interpretation}
\author {
    Shiao Xie\textsuperscript{\rm 1}\equalcontrib,
    Siyu Chen\textsuperscript{\rm 1}\equalcontrib,
    Jianwei Lv\textsuperscript{\rm 1},
    Bo Yuan\textsuperscript{\rm 1},
    Yujin Wang\textsuperscript{\rm 1}\corresponding,
    Xiandong Li\textsuperscript{\rm 1}\corresponding
}
\affiliations {
    % Affiliations
    \textsuperscript{\rm 1}Baidu Inc., China
    % firstAuthor@affiliation1.com, secondAuthor@affilation2.com, thirdAuthor@affiliation1.com
}
% \author{}
% \affiliations{}
% \author{
%     %Authors
%     % All authors must be in the same font size and format.
%     Written by AAAI Press Staff\textsuperscript{\rm 1}\thanks{With help from the AAAI Publications Committee.}\\
%     AAAI Style Contributions by Peter Patel Schneider,
%     Sunil Issar,\\
%     J. Scott Penberthy,
%     George Ferguson,
%     Hans Guesgen,
%     Francisco Cruz\equalcontrib\corresponding,
%     Marc Pujol-Gonzalez\equalcontrib\corresponding
% }
% \affiliations{
%     %Afiliations
%     \textsuperscript{\rm 1}Association for the Advancement of Artificial Intelligence\\
%     % If you have multiple authors and multiple affiliations
%     % use superscripts in text and roman font to identify them.
%     % For example,

%     % Sunil Issar\textsuperscript{\rm 2},
%     % J. Scott Penberthy\textsuperscript{\rm 3},
%     % George Ferguson\textsuperscript{\rm 4}\corresponding,
%     % Hans Guesgen\textsuperscript{\rm 5}
%     % Note that the comma should be placed after the superscript

%     1101 Pennsylvania Ave, NW Suite 300\\
%     Washington, DC 20004 USA\\
%     % email address must be in roman text type, not monospace or sans serif
%     proceedings-questions@aaai.org
% %
% % See more examples next
% }

%Example, Single Author, ->> remove \iffalse,\fi and place them surrounding AAAI title to use it
% \iffalse
% \title{My Publication Title --- Single Author}
% \author {
%     Author Name
% }
% \affiliations{
%     Affiliation\\
%     Affiliation Line 2\\
%     name@example.com
% }
% \fi

% \iffalse
% %Example, Multiple Authors, ->> remove \iffalse,\fi and place them surrounding AAAI title to use it
% \title{My Publication Title --- Multiple Authors}
% \fi

\nocopyright 
\usepackage{xr}
\externaldocument{Appendix}
\begin{document}

\maketitle

\begin{abstract}
% % 先抛出需求
% With the growing accessibility of medical reports through online healthcare platforms, patients increasingly seek personalized explanations of their reports, yet existing medical vision-language tasks do not adequately address this need. 
% % 给出任务的定义
% To bridge this gap, we introduce \emph{Patient-oriented Medical Report Interpretation} (PMRI), a novel open-ended multimodal generation task that requires models to explain medical reports according to a user's query and dialogue history in accurate and accessible language. 
% % 和现有任务的差异从而突出两个核心特点
% Unlike conventional medical report generation, PMRI simultaneously requires evidence-grounded medical factuality and context-dependent patient communication. 
% % 抛出任务的难点和挑战
% These objectives differ fundamentally in their verifiability, yet remain tightly coupled owing to the gap between clinical and lay registers, making them difficult to optimize under existing supervised fine-tuning and reinforcement learning paradigms.
% % 解决这个挑战我们的工作
% To address this challenge, we propose G-CARL, a grounded checklist-aligned reinforcement learning framework that integrates retrieval-grounded claim verification with case-specific checklist rewards. G-CARL verifies atomic medical claims through multi-source retrieval to supervise factuality, while using instance-specific weighted checklists to specify what a response should cover, enabling structured supervision for demand satisfaction and communication quality without constraining response diversity.
% % benchmark的建立
% We further construct \textsc{MMedReport}, a real-world PMRI benchmark together with a clinician-designed three-dimensional evaluation protocol. Extensive experiments demonstrate that G-CARL consistently outperforms existing post-training baselines in overall quality, claim-level precision, and checklist coverage. Pairwise clinician preference evaluation further confirms that G-CARL produces interpretations that are more accurate and better aligned with patient needs.

% while remaining faithful to report evidence and communicating in
% Clinical assertions can be adjudicated against the report and external clinical resources as objective facts, whereas need satisfaction and interpretive quality remain highly context-dependent. 
% Patient-oriented Medical Report Interpretation (PMRI) requires a multimodal model to interpret a medical report jointly with the user query and dialogue history, remaining faithful to the clinical knowledge in the report while conveying medically consequential content of user concern in lay register. 
% Clinical assertions can be adjudicated against the report and external clinical resources as objective facts, whereas need satisfaction and interpretive quality remain highly context-dependent. The two objectives therefore differ in verifiability yet remain coupled, owing to the cross-domain gap between clinical and lay registers. Existing SFT and RL methods handle them under a single conflated objective, failing to jointly guarantee medical correctness and adequate response to patient concerns. To address this limitation, we propose \textbf{G-CARL}, a checklist-aligned RL framework that integrates retrieval-grounded assertion verification with case-specific checklist rewards. G-CARL extracts atomic assertions and assesses their evidential support through multi-source retrieval, while generating a weighted, instance-level checklist for each case that prescribes only the content a response should cover without mandating reference wording, thereby supplying complementary reward signals for verifiable and context-dependent objectives. We further introduce MMedReport, a real-world multimodal PMRI benchmark with a clinician-designed three-dimensional evaluation protocol. On Qwen3-VL and InternVL3, G-CARL consistently surpasses SFT and RL baselines on quality score, assertion-level precision, and checklist coverage. Pairwise preference evaluation further confirms that G-CARL produces interpretations that are more accurate and more closely aligned with patient needs.

% Unlike conventional medical report generation, PMRI is simultaneously patient-centered and evidence-bounded: high-quality responses must address user-specific concerns in lay language while remaining strictly faithful to objective medical evidence. This dual nature gives rise to objectives with heterogeneous verifiability, making reward design particularly challenging for reinforcement learning. 
% To address this challenge, we propose \textbf{G-CARL}, a grounded and checklist-aligned reinforcement learning framework that decomposes reward supervision according to objective-specific verifiability. Specifically, medical factuality is optimized through retrieval-grounded claim verification, while demand satisfaction and expression quality are guided by dynamic checklists derived from physician references. 
% We further construct \textsc{MMedReport}, a real-world PMRI benchmark with a clinician-authored three-dimensional evaluation protocol. Experiments demonstrate that G-CARL consistently outperforms supervised fine-tuning and standard reinforcement learning baselines across all evaluation dimensions, while pairwise clinician preference studies further corroborate the effectiveness of our reward design.
% 先讲述什么是pmri任务
% Patient-oriented medical report interpretation (PMRI) requires multimodal models to explain uploaded medical reports in light of a user’s question and prior dialogue, while remaining faithful to report evidence and communicating in accessible language. 


% 突出不同维度的验证目标不同
% These objectives differ in their verifiability: medical claims can be verified against report content and external clinical resources, whereas demand satisfaction and communication quality depend on the specific consultation context.
% 提出我们的方法
% To address this challenge, we introduce \textbf{G-CARL}, a grounded checklist-aligned reinforcement learning framework that decomposes reward supervision into a \emph{Retrieval-Grounded Claim Reward} for medical factuality and a \emph{Case-Specific Checklist Reward} for context-dependent objectives.

% To address this challenge, we introduce \textbf{G-CARL}, a grounded checklist-aligned reinforcement learning framework that equiped with retrieval-grounded for medical claims 准确性 and case-specific checklist reward for回答的是否coverage.

% 具体的两个贡献是怎么做的
% Specifically, G-CARL extracts atomic claims and evaluates their report relevance and evidence support using multi-source retrieval. 
% it generated=sweighted, instance-specific checklists that specify \emph{what} a response should cover without enforcing a single reference wording. 
% % 提出了一个新的数据集以及评估指标
% We further construct \textsc{MMedReport}, a real-world multimodal PMRI benchmark and a clinician-designed three-dimensional evaluation protocol. 
% Across both Qwen3-VL and InternVL3 backbones, G-CARL consistently outperforms supervised fine-tuning and reinforcement learning baselines in clinician-defined quality scores, claim-level precision, and checklist coverage. 
% Pairwise clinician preference studies further confirm that G-CARL produces more accurate and patient-oriented interpretations.


% Patient-oriented Medical Report Interpretation (PMRI) requires a multimodal model to interpret a medical report in conjunction with a query and dialogue history, preserving fidelity to the clinical knowledge documented in the report while conveying, in lay register, the medically consequential content that addresses the user's concerns. Accordingly, clinical assertions can be adjudicated against the report and external clinical resources, whereas informational adequacy and interpretive quality remain contingent on the consultation context. The two classes of objectives therefore differ in verifiability, yet remain coupled through the cross-domain knowledge gap that both must traverse. Existing supervised fine-tuning (SFT) and reinforcement learning (RL) methods subsume them under a single undifferentiated objective and consequently struggle to jointly ensure factual correctness and an adequate response to patient concerns. We propose G-CARL, a checklist-aligned RL framework that integrates retrieval-grounded assertion verification with case-specific checklist rewards. It extracts atomic assertions and scores them by evidential support under multi-source retrieval, while instantiating a per-case weighted checklist that prescribes the requisite content without mandating reference wording, thereby supplying complementary rewards for verifiable and context-dependent objectives, respectively. We further construct MMedReport, a real-world multimodal PMRI benchmark, together with a three-dimensional evaluation protocol designed by practicing clinicians. Instantiated on two multimodal large language models (Qwen3-VL and InternVL3), G-CARL consistently surpasses SFT and RL baselines in overall quality, assertion-level precision, and checklist coverage. Pairwise preference evaluation further corroborates that G-CARL is preferred for both accuracy and alignment with the concerns expressed in the consultation.
% 首先提出需求和现状
Personalized interpretation of medical reports has emerged as an increasingly important need among patients. Addressing this need requires both evidence-grounded medical factuality and context-dependent patient communication, yet existing medical vision-language tasks do not adequately capture these dual requirements.
% 给出任务的定义
To bridge this gap, we introduce \emph{Patient-oriented Medical Report Interpretation} (PMRI), a novel open-ended multimodal generation task that requires models to explain medical reports in accurate and accessible language based on a user's query and dialogue history.
% 抛出任务的难点和挑战
% These objectives differ fundamentally in their verifiability, yet remain tightly coupled owing to the gap between clinical and lay registers, making them difficult to optimize under existing supervised fine-tuning and reinforcement learning paradigms.
These two objectives differ fundamentally in their verifiability, yet remain tightly coupled, making them difficult to optimize jointly under conventional supervised fine-tuning and holistic reinforcement learning paradigms.
% 解决这个挑战我们的工作
% To address this challenge, we propose G-CARL, a grounded, checklist-aligned reinforcement learning framework that integrates retrieval-grounded claim verification with instance-specific checklist rewards. G-CARL verifies atomic medical claims through multi-source retrieval to supervise factuality, while using instance-specific, context-aware weighted checklists to specify what a response should cover, enabling structured supervision for demand satisfaction and communication quality without constraining response diversity.
To address this challenge, we propose G-CARL, a grounded, checklist-aligned reinforcement learning framework that combines multi-source retrieval for atomic claim verification with context-aware, instance-specific weighted checklists for response coverage, providing structured supervision for factuality, user-demand satisfaction, and expression quality without constraining response diversity.
% benchmark的建立
We further construct \textsc{MMedReport}, a real-world PMRI benchmark, along with a clinician-designed three-dimensional evaluation protocol. Extensive experiments demonstrate that G-CARL consistently outperforms existing post-training baselines in overall quality, claim-level precision, and checklist recall. Pairwise preference evaluation by clinicians further confirms that G-CARL produces interpretations that are more accurate and better aligned with patient needs.
\end{abstract}

% Patient-oriented medical report interpretation (PMRI) requires multimodal models to explain uploaded medical reports in light of a user's query and prior dialogue while remaining faithful to report evidence and communicating in accessible language. Unlike conventional medical report generation, PMRI is simultaneously patient-centered and evidence-bounded. These heterogeneous objectives differ in their verifiability: medical claims can be verified against report content and external clinical resources, whereas demand satisfaction and communication quality depend on the specific consultation context.

% We introduce \textbf{G-CARL}, a verifiability-aware reinforcement learning framework that decomposes reward supervision accordingly. For medical factuality, G-CARL extracts atomic claims and evaluates their report relevance and evidence support through multi-source retrieval. For context-dependent objectives, it converts physician-authored reference interpretations into weighted, instance-specific checklists that specify what a response should cover without enforcing a single reference wording.

% We evaluate G-CARL on \textsc{MMedReport}, a multimodal benchmark containing report images, dialogue histories, user queries, and physician-authored reference interpretations. Across Qwen3-VL and InternVL3 backbones, G-CARL consistently improves clinician-defined quality scores and case-specific checklist coverage over supervised fine-tuning and representative reinforcement learning baselines. Additional clinician preference studies further demonstrate that G-CARL produces more accurate, relevant, and patient-friendly medical interpretations.
% Uncomment the following to link to your code, datasets, an extended version or similar.
% You must keep this block between (not within) the abstract and the main body of the paper.
% Make sure that you do not de-anonymize yourself with these links.
% \begin{links}
%     \link{Code}{https://aaai.org/example/code}
%     \link{Datasets}{https://aaai.org/example/datasets}
%     \link{Extended version}{https://aaai.org/example/extended-version}
% \end{links}
\section{Introduction}
% 第一段重点讲述一下pmri这个任务，描述下他的临床价值
% Medical reports are central to clinical decision making, but they are difficult for patients to interpret. 
Written within a professional medical context, medical reports commonly present abnormal values and descriptive findings without explaining their implications for non-expert readers. 
This gap becomes particularly salient in online healthcare scenarios, where patients upload one or more report images and ask open-ended questions shaped by their personal concerns. 
To bridge this gap, we introduce \emph{Patient-oriented Medical Report Interpretation} (PMRI), a novel open-ended multimodal generation task.
% PMRI is designed to bridge this gap. 
Beyond simply recognizing report findings, PMRI must align professional medical knowledge with patient-centered communication by transforming report evidence into accessible explanations, addressing user-specific concerns, and offering clinically cautious guidance.

\begin{figure}[t]
    \centering
\includegraphics[width=0.99\linewidth]{intro2.pdf}
    \caption{
    % Comparison of PMRI (\textcolor{myorange}{\textit{\textbf{Task C}}}) with medical VQA (\textit{\textbf{Task A}}) and conventional report generation (\textit{\textbf{Task B}}). While medical VQA answers a single query with a deterministic answer and report generation describes image findings in a short form, PMRI integrates multimodal medical reports with long patient--doctor dialogue contexts to generate long-form, evidence-bounded, and patient-facing explanations.
    Comparison of PMRI (\textcolor{myorange}{\textit{\textbf{Task C}}}) with medical VQA (\textit{\textbf{Task A}}) and conventional report generation (\textit{\textbf{Task B}}). Unlike the short, deterministic outputs of Tasks A and B, PMRI produces long-form, patient-facing explanations grounded in multimodal reports and extended patient–doctor dialogue.
    }
    \vspace{-16pt}
    \label{fig:task_comparison}
\end{figure}

% 第二段需要系统的说明这个任务和之前任务的差异，最重要的两点是evidence-bounded以及patient-facing的特性，然后引出我们建模的三个评估维度：医疗准确性、需求满足度以及表达效果
As illustrated in Fig.~\ref{fig:task_comparison}, PMRI differs from conventional medical visual question answering (VQA)~\cite{lu2025bridging,jiang2025knowing} and report generation~\cite{lin2026march,li2025joint} in two important ways. First, it is an \emph{evidence-bounded} generation task whose responses must be grounded not only in the reports but also in reliable clinical knowledge and medically coherent diagnostic reasoning, particularly when explaining abnormalities or suggesting follow-up actions. This requirement is essential for patient safety because unsupported interpretations and inappropriate recommendations may mislead patients or delay necessary care. Second, PMRI is a \emph{patient-facing} communication task that must go beyond clinical correctness to address patients' individual information needs while clearly explaining potential abnormalities and risks in an emotionally adaptive and reassuring manner. Accordingly, we model PMRI quality along three core dimensions: medical accuracy, demand satisfaction, and expression quality.
% making it a user-conditioned, evidence-bounded, and safety-critical generation task.

% Second, it is a \emph{patient-facing} communication task, where clinical correctness alone is insufficient. Since users often seek reassurance, clarification, or practical next steps, a high-quality response must satisfy the user's information need while communicating uncertainty and medical implications in a clear and responsible manner. 

% A single reference interpretation cannot cover this response space, making reference-based supervision alone insufficient for learning diverse personalized explanations that remain clinically grounded.
% the dual nature of the task. On the one hand, medical factuality should be grounded in the uploaded report and reliable medical knowledge, since unsupported claims may lead to unsafe advice. On the other hand, user-need satisfaction and expression quality are context-dependent, because a useful response must address the user's specific concern in an appropriate manner. PMRI therefore requires rewards that can support both evidence-grounded medical correctness and patient-centered communication quality.


% 这一段主要描述下一般的sft不行的原因，sft单纯模仿正确答案，但是答案不是唯一的
Given the nature of PMRI, a straightforward strategy is to collect large-scale physician-written interpretations and train multimodal models with supervised fine-tuning (SFT)~\cite{chen2024sharegpt4v,rotstein2024fusecap}. However, SFT can overfit to the specific wording of the reference interpretation and encourage imitation of particular answers rather than learning the underlying principles. This is especially limiting in PMRI, where multiple responses may be clinically acceptable for the same report and user query as long as they remain faithful to the evidence and address the user's concern. Physician references therefore provide useful guidance on what a response should cover, rather than fully defining the quality space of PMRI outputs.


Reinforcement learning (RL) with reward-based post-training, such as Group Relative Policy Optimization (GRPO)~\cite{grpo}, offers a promising alternative for improving large vision-language models beyond reference imitation~\cite{xing2025caprl}. 
However, applying RL to PMRI requires rewards that reflect heterogeneous verifiability of different objectives. Medical factuality can be externally verified against the uploaded report and clinical knowledge, whereas demand satisfaction and expression quality depend more on the patient's concern and the consultation context. This motivates reward signals that can distinguish objective-specific errors rather than collapsing the entire response into a single score. 


Existing reward designs do not fully meet this requirement. Holistic MLLM-as-a-Judge scoring~\cite{zheng2023judging,chen2024mllm} compresses an entire interpretation into a single reward, making the supervision coarse and highly dependent on the judge model's medical knowledge. As a result, localized hallucinations may be overlooked when the overall response appears fluent and plausible, despite the fact that a single unsupported medical claim can fundamentally mislead patients. Recent work has introduced rubric-based rewards to provide more structured supervision~\cite{arora2025healthbench,gunjal2025rubrics}. However, static rubrics remain insufficient for PMRI because evaluation priorities vary substantially across cases. 
PMRI errors often stem not from entirely incorrect responses, but from omitting case-critical information or emphasizing secondary details while overlooking the most important clinical recommendations and user concerns. Since static rubrics are designed to capture generic response quality, they are often insensitive to these case-specific omissions and misplaced emphases.

% as fixed criteria are difficult to align with the specific report and user question. A patient asking about disease risk, follow-up necessity, or the meaning of an abnormal indicator may require different evaluation focuses, and query-agnostic rubrics may miss the information need and report-specific errors that matter most in that consultation.


To address these challenges, we propose G-CARL, a reinforcement learning framework with retrieval-grounded and checklist-guided rewards. 
G-CARL assigns reward mechanisms according to the verifiability boundary of each objective. For externally verifiable medical factuality, it decomposes each response into atomic medical claims, retrieves supporting evidence from the uploaded report and a multi-source medical datastore, and evaluates each claim for factual support and contextual relevance. This claim-level reward provides localized supervision for sparse factual errors and discourages unsupported elaboration.
For context-dependent objectives such as demand satisfaction and expression quality, G-CARL constructs instance-specific weighted checklists through MLLM generation followed by clinician refinement. Each checklist item is assigned an automatically generated weight, allowing the reward to emphasize the aspects most relevant to the current report and user question. This yields an explicit checklist score that provides transparent supervision for whether the response addresses the user's concern appropriately.
In summary, our contributions are as follows:
\begin{itemize}

\item We formulate PMRI as an evidence-grounded and patient-facing multimodal generation task, and propose G-CARL, a reinforcement learning framework that decomposes reward supervision according to the heterogeneous verifiability of different objectives.

\item We propose a retrieval-grounded claim reward that provides fine-grained supervision for medical factuality through atomic claim verification, and a case-specific checklist reward that explicitly supervises demand satisfaction and expression quality without relying on a single reference response.

\item We construct \textsc{MMedReport}, a real-world PMRI benchmark with clinician-designed evaluation protocols. Extensive experiments across multiple LVLM backbones demonstrate that G-CARL consistently outperforms supervised and reinforcement learning baselines, with the gains further validated by clinician preference and user comprehension studies.
\end{itemize}

% \textbf{\textcolor{deepred}{

% \begin{itemize}

% \item We formulate patient-oriented medical report interpretation as a novel evidence-bounded and patient-facing multimodal generation task, and propose G-CARL, a reinforcement learning framework that explicitly models heterogeneous supervision according to the verifiability of different objectives.

% \item We propose a retrieval-grounded claim verification reward that decomposes long-form interpretations into atomic medical claims and verifies each claim against report evidence and a multi-source medical knowledge datastore, providing fine-grained factual supervision for reinforcement learning.

% \item We introduce a weighted dynamic checklist reward that converts physician-authored reference interpretations into instance-specific evaluation criteria with adaptive importance weights, enabling structured supervision for demand satisfaction and expression quality.

% \item We construct \textsc{MMedReport}, a real-world PMRI benchmark with clinician-authored evaluation protocols, and conduct experiments across multiple LVLM backbones. Both automatic evaluation and clinician preference studies demonstrate the effectiveness of G-CARL over strong supervised and reinforcement learning baselines.

% \end{itemize}


% You can remove the copyright notice and ensure that your names aren't shown by including \texttt{submission} option when loading the \texttt{aaai2027} package:
\begin{figure*}[t]
    \centering
    \includegraphics[width=0.95\linewidth]{method2.pdf}
    \vspace{-8pt}
    \caption{
    Overview of G-CARL. Given report images, user queries, and dialogue history, G-CARL samples $G$ responses from the old policy and optimizes them with a multi-objective reward function consisting of retrieval-grounded claim reward $R_{\mathrm{fact}}$, case-specific checklist reward $R_{\mathrm{check}}$, and structured reasoning format reward $R_{\mathrm{format}}$.
   }
    \label{fig:G-CARL_architecture}
\end{figure*}

% Overview of G-CARL. Given input report images, a user query, and dialogue history, we first sample $G$ candidate responses $\{o_i\}_{i=1}^{G}$ from the old policy. 
% Each response $o$ is optimized using a multi-objective reward function comprising three components: 
%    (1) a retrieval-grounded claim reward $R_{\mathrm{fact}}$ for externally verifiable medical factuality, 
%    (2) a case-specific checklist reward $R_{\mathrm{check}}$ for demand satisfaction and expression quality, and (3) a format reward $R_{\mathrm{format}}$ that encourages structured diagnostic reasoning within the \texttt{<think>}\texttt{</think>} block.
\section{Related Works}
\noindent \textbf{Medical Report Generation}. 
Medical report generation has been extensively studied where models generate diagnostic reports from multimodal images such as X-rays~\cite{Li_2023_CVPR,Liu_2025_CVPR}. Recent methods have improved report quality through multimodal feature alignment~\cite{jin2024promptmrg,li2025joint}, clinically grounded visual representations~\cite{varol2025vision}, and reinforcement learning-based optimization~\cite{wang2026beyond}. 
MedVAG~\cite{varol2025vision} introduces clinically aware visual grounding, while HiMed-3B~\cite{wang2026beyond} explores RL-based alignment for medical text generation. MedRepBench~\cite{shang2025medrepbench} further promotes faithful report generation through field-level evaluation of structured clinical findings. In contrast to these imaging-to-report tasks, PMRI focuses on patient-facing interpretation of existing structured reports under patient--doctor dialogue contexts, requiring models to address user-specific concerns while providing accurate and understandable explanations. 


% \section{Related Works}
% \noindent \textbf{Medical Report Generation}.
% Medical report generation has been extensively studied where models generate diagnostic reports from multimodal images such as X-rays.
% Earlier approaches established knowledge- and memory-driven paradigms, including KERP, R2Gen, and R2GenCMN, which leverage structured clinical knowledge and relational or cross-modal memories to improve long-form radiology report generation~\cite{li2019kerp,chen2020r2gen,chen2021r2gencmn}.
% Recent methods have improved report quality through multimodal feature alignment~\cite{jin2024promptmrg,li2025joint}, clinically grounded visual representations~\cite{varol2025vision}, and reinforcement learning-based optimization~\cite{wang2026beyond}.
% MedVAG~\cite{varol2025vision} introduces clinically aware visual grounding, while HiMed-3B~\cite{wang2026beyond} explores RL-based alignment for medical text generation. MedRepBench~\cite{shang2025medrepbench} further promotes faithful report generation through field-level evaluation of structured clinical findings.
% LLM-based systems such as R2GenGPT and RaDialog further align radiographic representations with pretrained language models, with RaDialog additionally supporting interactive report correction and question answering~\cite{wang2023r2gengpt,pellegrini2023radialog}. Region-guided generation improves anatomical transparency and clinician controllability by describing localized anatomical regions before composing the final report~\cite{tanida2023region}.
% In contrast to these clinician-oriented tasks, PMRI focuses on patient-facing interpretation of existing medical reports under multi-turn dialogue contexts, requiring models to address user-specific concerns while providing accurate and understandable explanations.
% Closest to PMRI, patient-centered studies have explored translating radiology reports into lay language and accessible multimedia explanations~\cite{jeblick2023simplified,luo2024rexplain}; however, they generally focus on one-shot reformulation of a single radiology report rather than personalized interpretation across heterogeneous reports and multi-turn user interactions.

% \vspace{2pt}
% \noindent \textbf{Reinforcement Learning for Medical VLMs}.
% Reinforcement learning (RL) has recently been adopted to improve reasoning and reliability in medical VLMs.
% Earlier report-specific RL methods optimized clinical entity completeness, inferential consistency, semantic-graph agreement, and cross-modal alignment instead of relying solely on lexical-overlap objectives~\cite{miura2021factual,delbrouck2022semantic,qin2022reinforced}.
% MedVLM-R1~\cite{pan2025medvlm} uses a GRPO-based framework to elicit explicit reasoning paths for radiology VQA, while Med-R1~\cite{lai2026med} designs preference signals that align visual perception, intermediate reasoning, and final answers.
% Beyond short-form QA, MediX-R1~\cite{mullappilly2026medix} extends multimodal medical RL to open-ended responses through LLM-based multi-objective rewards.
% RL has also been explored for report-centric tasks.
% RadVLM-GRPO~\cite{gundersen2026radvlm} applies clinically grounded rewards to chest X-ray report generation and visual grounding, showing that RL can complement strong supervised fine-tuning.
% Complementary report-centric alignment methods include multi-objective preference optimization and online iterative self-alignment, which balance clinical accuracy, completeness, and linguistic quality~\cite{xiao2025mpo,xiao2025oisa}. EditGRPO and OraPO further introduce post-rollout corrections or oracle-guided preference supervision to improve exploration efficiency and factuality~\cite{zhang2025editgrpo,chen2026orapo}.


\vspace{1pt}
\noindent \textbf{Reinforcement Learning for Medical VLMs}. 
RL has recently been adopted to improve reasoning and reliability in medical VLMs~\cite{jing2026reason,zhou-etal-2026-enhancing}. 
MedVLM-R1~\cite{pan2025medvlm} uses a GRPO-based framework to elicit explicit reasoning paths for radiology VQA, while Med-R1~\cite{lai2026med} designs preference signals that align visual perception, intermediate reasoning, and final answers. 
Beyond short-form QA, MediX-R1~\cite{mullappilly2026medix} extends multimodal medical RL to open-ended responses through LLM-based multi-objective rewards. 
RL has also been explored for report-centric tasks. 
RadVLM-GRPO~\cite{gundersen2026radvlm} applies clinically grounded rewards to chest X-ray report generation and visual grounding, showing that RL can complement strong SFT.  




\section{Methods}
% \subsection{Problem Definition}
% We formalize patient-oriented medical report interpretation as a multimodal conditional generation problem. Each input instance $x$ consists of a patient query $q$, a sequence of dialogue history turns $h$, and one or more uploaded medical report images $I$. The goal is to learn a vision-language policy $\pi_\theta$ that produces an interpretation $y$ which (i) remains faithful to the medical evidence available in $I$ and reliable external knowledge, (ii) directly addresses the user's specific concern, and (iii) communicates the explanation in clinically safe and patient-accessible language.

\subsection{Overall Architecture}
 % illustrates the overall framework of G-CARL. 
Built upon GRPO, G-CARL optimizes the policy model with three reward signals, as illustrated in Fig.~\ref{fig:G-CARL_architecture}. Given uploaded medical report images $I$, the dialogue history $h$, and a user query $q$, the policy model first generates candidate responses. The retrieval-grounded branch then supervises medical factuality by verifying report-grounded medical claims against evidence retrieved from a multi-source medical datastore, producing the factuality reward $R_{\mathrm{fact}}$ (\textbf{\textcolor{deepred}{Sec.~\ref{retrieval_grounded}}}). In parallel, the case-specific checklist branch optimizes demand satisfaction and expression quality by constructing a weighted checklist through MLLM generation followed by clinician-guided refinement, and evaluating the checklist coverage of the generated response to produce the checklist reward $R_{\mathrm{check}}$ (\textbf{\textcolor{deepred}{Sec.~\ref{case_specific_check}}}). Together with the format reward $R_{\mathrm{format}}$, these components are integrated into a multi-objective optimization task (\textbf{\textcolor{deepred}{Sec.~\ref{multi_obj_reward}}}).

%  which are then evaluated through two parallel reward branches. 

\subsection{Retrieval-Grounded Claim Reward}
\label{retrieval_grounded}
Medical factuality in PMRI is inherently claim-level, as a single response often contains multiple heterogeneous medical claims. Holistic response-level rewards cannot localize factual errors and may overestimate fluent but unsupported generations. Moreover, many clinical interpretations, causal attributions, and recommendations cannot be verified from the report alone, requiring factual verification grounded in authoritative external medical knowledge. 

To address these challenges, inspired by~\cite{factscore,veriscore}, we propose a retrieval-grounded claim  reward. 
As illustrated in Fig.~\ref{fig:G-CARL_architecture}~(a), each response is first decomposed into atomic medical claims, which are then verified against evidence retrieved from a multi-source medical datastore. 
Specifically, each claim is evaluated from two complementary perspectives: whether it is \emph{RELEVANT} to the uploaded medical report and whether it is \emph{SUPPORTED} by the retrieved evidence.
This dual-binary design encourages the generation of statements that are both clinically grounded and factually correct. We describe each step in detail below.

\vspace{2pt}
\noindent \textbf{Claim Extraction.}
Following~\cite{veriscore}, we perform claim extraction at the sentence level. Given a response $o$, we first segment its answer span into sentences $\{s_1, \dots, s_n\}$. For each sentence $s$, an extractor $\mathcal{G}$, conditioned on the full response $o$ and user query $q$, produces a set of verifiable and decontextualized claims as follows:
\begin{equation}
\{(c, \tilde{q}), \dots\}
=
\mathcal{G}(s \mid o, q).
\end{equation}
Alongside each claim $c$, the extractor jointly generates a retrieval query $\tilde{q}$. 
After claim deduplication, the final claim set for response $o$ is defined as:
\begin{equation}
\mathcal{C}(o)=\{(c_i,\tilde{q}_i)\}_{i=1}^{N}.
\end{equation}


\vspace{2pt}
\noindent \textbf{Evidence Retrieval.}
To support trustworthy evidence-grounded factual verification in the medical domain, we construct a large-scale medical datastore $\mathcal{D}$ by integrating authoritative medical resources, including drug labeling, medical textbooks, and clinical practice guidelines:
\begin{equation}
\mathcal{D} = \mathcal{D}_{\text{drug}} \cup \mathcal{D}_{\text{book}} \cup \mathcal{D}_{\text{guide}},
\end{equation}
where $\mathcal{D}_{\text{drug}}$ contains approximately 20,000 drug instruction entries, $\mathcal{D}_{\text{book}}$ comprises 18,200 textbook passages covering foundational medical knowledge, and $\mathcal{D}_{\text{guide}}$ consists of 16,000 clinical guideline passages. All resources are preprocessed into semantically coherent chunks.

% and indexed by a dense retriever $f$.
Given a claim-query pair $(c, \tilde{q})$, the retrieval query $\tilde{q}$ combines the topical intent of the user query with the core medical concepts of the corresponding claim, producing a concise set of retrieval keywords aligned with the current clinical context. We then perform parallel retrieval across all knowledge sources using $\tilde{q}$ and aggregate the top-$k$ most relevant evidence chunks from each source:
% To accommodate the diverse medical knowledge required for accurate report interpretation, we organize $\mathcal{D}$ into specialized sub-collections with distinct functional roles:
% \begin{itemize}[leftmargin=*,itemsep=0pt,topsep=0pt]
%     \item $\mathcal{D}_{\text{book}}$ provides foundational medical facts, disease mechanisms, and pathophysiological explanations;
%     \item $\mathcal{D}_{\text{guide}}$ offers evidence-based clinical guidelines for diagnosis, treatment, and disease management;
%     \item $\mathcal{D}_{\text{taboo}}$ specifies contraindications, high-risk procedures, and prohibited actions for patient safety.
% \end{itemize}
% therefore, we perform parallel retrieval across all sources and aggregate the top-$k$ most relevant chunks from each source, yielding the following formulation as:
\begin{equation}
\mathcal{E} = \bigcup \mathrm{Top}_k\big(f( \tilde{q}, \mathcal{D}_i)\big), \quad i \in \{\text{book}, \text{guide}, \text{drug}\},
\end{equation}
where $f(\tilde{q}, d)$ denotes the relevance score between query and document chunk. 
The retrieved evidence set $\mathcal{E}$ is then concatenated with source identifiers and truncated to a fixed token budget to form the evidence context for verifying claim. 
% This multi-source retrieval strategy enables different categories of claims to be grounded in the most appropriate bodies of medical knowledge.

\vspace{2pt}
\noindent \textbf{Dual Binary Verification.}
To verify each extracted claim, a multimodal verifier $\mathcal{V}$ takes the claim $c$, the retrieved evidence $\mathcal{E}$, the user query $q$, and the uploaded report image $I$ as input, and predicts two binary verdicts as:
\begin{equation}
\begin{aligned}
v &= \mathcal{V}(c, \mathcal{E}),
& v &\in \{\textsc{Supported}, \textsc{UnSupported}\},\\
u &= \mathcal{V}(c, I, h, q),
& u &\in \{\textsc{Relevant}, \textsc{Irrelevant}\}, 
\end{aligned}
\end{equation}
where $v$ assesses whether the claim is factually supported, while $u$ assesses whether the claim is contextually relevant to the uploaded reports. 

For factuality verification, the verifier follows an evidence-first, contradiction-oriented protocol. A claim is judged as \textsc{Supported} if it is directly supported by the retrieved evidence or remains consistent with it without contradiction. When explicit evidence is unavailable, the claim is still considered \textsc{Supported} if it reflects established medical consensus and does not conflict with domain knowledge. Otherwise, the claim is judged as \textsc{Unsupported}.

%  if it is contradicted by the retrieved evidence or cannot be reliably substantiated by authoritative medical knowledge

For relevance verification, a claim is judged as \textsc{Relevant} if it is grounded in the uploaded report, directly addresses the user query, or provides clinically necessary context for the current case. Otherwise, it is judged as \textsc{Irrelevant}, even if medically correct, when it introduces generic or unsupported information unrelated to the report findings or the user's concern. This dual-verdict design prevents the verifier from penalizing valid common-sense medical statements that lack explicit retrieved evidence, while discouraging reward hacking through factually correct but irrelevant elaborations. 

Then, a claim is considered valid only if it is both factually supported and contextually relevant: 
\begin{equation}
z
=
\mathbf{1}[v=\textsc{Supported}]
\cdot
\mathbf{1}[u=\textsc{Relevant}]
\in
\{0,1\}.
\end{equation}
% Let $N$ denotes the total number of extracted claims and
% $S=\sum_{i=1}^{N} z_i$ denote the number of valid claims. 
The valid-claim precision can be defined as:
\begin{equation}
P_{\mathrm{fact}}
=
\frac{1}{N}
\sum_{i=1}^{N} z_i.
\end{equation}
However, precision alone may favor overly short responses, as a response containing only a few valid claims can still achieve a high precision score. 
% To encourage sufficient informational coverage, we introduce a capped coverage term with a case-specific target $K$, 由~\ref{case_specific_check}中的generator自动生成， which specifies the expected number of valid medical claims in an informative response:
To encourage sufficient informational coverage, we introduce a coverage term with a case-specific target claim count $K$, where $K$ is automatically generated by the checklist generator described in Section~\ref{case_specific_check} and represents the expected number of valid medical claims in an informative response:
\begin{equation}
C_{\mathrm{fact}}
=
\min\!\left(
\frac{\sum_{i=1}^{N} z_i}{K}, 
1
\right).
\end{equation}
The final factual validity reward can be formulated as the harmonic mean of valid-claim precision and coverage: 
\begin{equation}
R_{\mathrm{fact}}
=
\frac{2P_{\mathrm{fact}}C_{\mathrm{fact}}}{P_{\mathrm{fact}}+C_{\mathrm{fact}}}.
\end{equation}
% where $\epsilon$ is a small constant for numerical stability.


\subsection{Case-Specific Checklist Reward}
\label{case_specific_check}
Unlike medical factuality, dimensions such as demand satisfaction and expression quality are inherently subjective and cannot be verified against referenceable knowledge. 
Moreover, these dimensions are highly instance-dependent: 
whether a response is considered complete, helpful, or well-expressed depends on the specific report, user concern, and dialogue context. 
Different criteria also vary substantially in clinical importance. Consequently, assigning a single holistic reward provides little guidance about which aspects of the response should be improved and fails to distinguish critical requirements from desirable but non-essential ones.

To address this challenge, we decompose subjective evaluation into a set of weighted checklist items. Each checklist item represents an independent evaluation criterion with an associated importance weight. A verifier then evaluates the candidate response against each checklist item independently, producing a fine-grained and controllable reward signal for reinforcement learning. The overall procedure is described as follows.


%  dynamically synthesized from physician-authored reference interpretations
% \subsection{Case-Specific Checklist Reward}
% \label{case_specific_check}
% Unlike medical factuality, dimensions such as demand satisfaction and expression quality are inherently subjective and cannot be directly verified against external knowledge sources. Moreover, these dimensions are highly instance-dependent: whether a response is considered complete, helpful, or well-expressed varies with the specific medical report, user concerns, and dialogue context. Meanwhile, the relative importance of different criteria can differ substantially across cases. Consequently, a single holistic reward provides limited guidance on which aspects of a response require improvement and fails to distinguish clinically critical requirements from desirable but non-essential qualities.To address this challenge, we decompose subjective evaluation into a set of weighted checklist items dynamically synthesized from physician-authored reference interpretations. Each checklist item corresponds to an independent evaluation criterion and is assigned an importance weight reflecting its relative significance. A verifier then assesses the candidate response against each checklist item separately, generating a fine-grained, interpretable, and controllable reward signal for reinforcement learning. The overall procedure is described as follows.

% To supervise these subjective yet clinically important dimensions, rather than relying on a single holistic reward score, we decompose evaluation into a set of explicit checklist items dynamically synthesized from physician-authored reference interpretations. A judge then evaluates the candidate response against these checklist items independently, yielding structured and interpretable supervision for RL optimization.

\vspace{2pt}
\noindent \textbf{Weighted Checklist Construction.}
As shown in Fig.~\ref{fig:G-CARL_architecture}~(b), to reduce manual annotation effort, an MLLM (\textit{e.g.}, Gemini 3.1 Pro) as generator $G$ first generates a draft instance-specific checklist conditioned on the dialogue history, user query, and uploaded medical reports. Professional clinicians then iteratively review, refine, and complete the draft to produce the final checklist:
\begin{equation}
\mathcal{T}
=
\{(t_i, w_i)\}_{i=1}^{m},
\end{equation}
where $t_i$ denotes a self-contained checklist item and $w_i$ its corresponding importance weight. 
Physician-guided checklist construction grounds the supervision signal in clinically appropriate expectations rather than generic evaluation templates, which prior work has shown to be essential for reliable expert-domain assessment~\cite{viswanathan2025checklists}.


The importance weight $w$ is determined according to the clinical significance of each checklist item. Specifically, each checklist item is assigned one of four importance levels, namely \textsc{Essential}, \textsc{Important}, \textsc{Optional}, or \textsc{Pitfall}, which are mapped to predefined integer-valued weights (e.g., 4--5, 2--3, 1--2, and $-2$--$-1$, respectively). This weighting scheme ensures that clinically critical criteria contribute more strongly to the reward while undesirable behaviors incur explicit penalties.
In particular, \textsc{Pitfall} items capture undesirable behaviors discouraged by the physicians, including ignoring the user's primary concern, generating overly technical explanations, or providing dismissive responses. 


\vspace{2pt}
\noindent \textbf{Explicit Aggregation.}
Given the checklist $\mathcal{T}$ and a candidate response $o$, the checklist verifier $\mathcal{V}$ evaluates each checklist item independently and outputs a binary satisfaction signal:
\begin{equation}
z_i = \mathcal{V}(t_i, o) \in \{0,1\}.
\end{equation}
For \textsc{Essential}, \textsc{Important}, and \textsc{Optional} items, $z_i=1$ indicates that the criterion is satisfied. For \textsc{Pitfall} items, $z_i=1$ indicates that the undesirable behavior is triggered. 
We denote the positive checklist item set as $\mathcal{P}$ and the \textsc{Pitfall} item set as $\mathcal{N}$, and compute the positive coverage and pitfall violation terms as: 
\begin{equation}
V_{\mathrm{pos}}
=
\frac{
\sum_{i\in\mathcal{P}} |w_i| z_i
}{
\sum_{i\in\mathcal{P}} |w_i|
},
\qquad
V_{\mathrm{neg}}
=
\frac{
\sum_{j\in\mathcal{N}} |w_j| z_j
}{
\sum_{i\in\mathcal{P}} |w_i|
}.
\end{equation}
The final checklist reward is defined as:
\begin{equation}
R_{\mathrm{check}}
=
\mathrm{clip}\!\left(
V_{\mathrm{pos}} - V_{\mathrm{neg}},
0,
1
\right).
\end{equation}
% where $\epsilon$ is a small constant for numerical stability. 
% We then compute the checklist reward by separately normalizing positive coverage and pitfall violations:
% We then adopt explicit checklist aggregation instead of a single continuous score, as independent evaluation of checklist items provides a more stable and interpretable supervision signal by localizing each decision to an individual criterion and reducing the conflation inherent in holistic judgments. The checklist reward is computed as the normalized weighted sum that can be formulated as:
% \begin{equation}
% R_{\mathrm{check}}
% =
% \frac{
% \sum_{i=1}^{m} w_i z_i
% }{
% \sum_{i=1}^{m} w_i
% }.
% \end{equation}
Unlike conventional reference-based training, our checklist does not encourage the policy to mimic a single physician-authored response. Instead, the reference is distilled into a set of clinically important evaluation criteria, specifying \emph{what} the response should cover rather than \emph{how} it should be written. This design preserves the diversity of valid responses while rewarding clinical completeness and user-oriented communication.
% Second, explicit weighting makes the reward composition transparent and controllable, so high-priority items such as \textsc{Essential} and \textsc{Pitfall} dominate the reward by construction, preventing superficial stylistic improvements from compensating for failures to satisfy the user's core informational needs.
% Finally, normalization by the total weight ensures reward comparability across instances containing different numbers of checklist items.

\subsection{Multi-objective Reward Function}
\label{multi_obj_reward}
Inspired by recent advances in structured reasoning~\cite{yu2025docthinker}, we introduce a format reward that encourages the model to follow a physician-oriented clinical reasoning scaffold during reinforcement learning. As illustrated in Fig.~\ref{fig:G-CARL_architecture}, the model is required to organize its intermediate reasoning within the \texttt{<think>}\texttt{</think>} block according to a four-step clinical reasoning workflow before generating the final patient-facing interpretation in the \texttt{<answer>}\texttt{</answer>} block. This reward does not directly supervise medical correctness; instead, it encourages a structured decomposition of report findings, supporting evidence, clinical reasoning, and response planning, thereby improving the coherence and organization of the final response.
Our reward function jointly optimizes medical factuality, demand satisfaction and expression quality, which can be defined as:
\begin{equation}
R_\text{total}
=
\lambda_{\text{fact}}R_\text{fact}
+
\lambda_{\text{check}}R_\text{check}
+
\lambda_{\text{format}}R_\text{format},
\end{equation}
where $\lambda_{\mathrm{fact}}$, $\lambda_{\mathrm{check}}$, and $\lambda_{\mathrm{format}}$ are weighting coefficients that balance the contributions of the reward components.


% This structured reasoning paradigm improves physician-aligned diagnostic coherence and enhances the interpretability of the reasoning process. 


\section{Experiments}

\subsection{Experimental Setup}
\noindent \textbf{Datasets.} 
% We evaluate G-CARL on \textsc{MMedReport}, a large-scale real-world multimodal PMRI benchmark collected from online healthcare consultations. The dataset contains 2,450 consultation instances, each consisting of the dialogue history, user query, uploaded medical report images, and a physician-authored reference checklist. To ensure data quality and patient privacy, all instances undergo quality control, de-identification, and manual verification. Detailed dataset statistics analyses are provided in the \textbf{Appendix~A}. Following the standard split, 2,200 instances are used for training and the remaining 250 instances for evaluation.
% we additionally evaluate the trained model on two external medical question-answering benchmarks, \textsc{CMB}~\cite{} and \textsc{Chinese Medical VQA}~\cite{}.
We evaluate G-CARL primarily on \textbf{(\textcolor{deepred}{1})} \textsc{MMedReport}, a real-world multimodal PMRI benchmark collected from online healthcare consultations. It contains 2,450 instances, each including dialogue history, a user query, uploaded medical report images, and clinician-verified reference annotations. All instances undergo quality control, de-identification, and manual verification to ensure data quality and patient privacy. Detailed dataset statistics are provided in \textbf{Appendix~A}. Following the standard split, 2,200 instances are used for training and 250 for evaluation. 
To assess broader medical capability beyond PMRI, we further evaluate the trained models on an external benchmark, \textbf{(\textcolor{deepred}{2})} \textsc{CMB}~\cite{wang2024cmb}, which evaluates both medical QA accuracy and the professionalism of open-ended clinical interpretation under an LLM-as-a-judge protocol.

% and \textbf{(\textcolor{deepred}{3})} Chinese SimpleVQA~\cite{gu2025see}.
% 额外数据集测试
% Because PMRI involves subjective communication dimensions, relying solely on unguided LLM judges is insufficient. We introduce a clinician-authored evaluation protocol. 
% This decouples the training reward (instance-specific) from the evaluation metric (global dimension-level criteria).

\vspace{1pt}
\noindent \textbf{Evaluation Protocol.} We evaluate PMRI responses using both subjective and objective metrics. (1) \emph{Subjective Evaluation}.
Responses are evaluated along three clinician-defined dimensions: \emph{Medical Accuracy} (Accuracy), \emph{Demand Satisfaction} (Satisfaction), and \emph{Expression Quality} (Expression). Each dimension is scored according to a detailed clinician-authored rubric ranging from $-2$ to $3$. We employ GPT-5.2~\cite{openai2025gpt52systemcard} as the judge by strictly following these evaluation criteria. 
The complete judging protocol is provided in the \textbf{Appendix B} to facilitate reproducibility.
(2) \emph{Objective Evaluation}.
In addition to holistic assessment, we report claim-level \emph{Precision} and checklist-level \emph{Recall}. Precision is defined as the ratio of supported claims to extracted claims, measuring factual correctness, while Recall is defined as the ratio of satisfied checklist items to the total checklist items, measuring coverage of case-specific requirements. 

\vspace{1pt}
\noindent \textbf{Implementation Details.}
Our policy models are initialized from the Qwen3-VL~\cite{bai2025qwen3vltechnicalreport} and InternVL3~\cite{zhu2025internvl3exploringadvancedtraining} series. G-CARL is trained for three epochs on 8 NVIDIA H100 GPUs with a batch size of 192 and a learning rate of $1\times10^{-5}$. 
We set the number of rollout responses to $G=8$ during training. 
The reward weights are configured as $\lambda_{\text{fact}}=0.4$, $\lambda_{\text{check}}=0.3$, and $\lambda_{\text{format}}=0.3$; detailed hyperparameter analysis is provided in \textbf{Appendix C}. 
The verifier $\mathcal{V}$ is initialized with Qwen3.5-35B-A3B, with implementation details reported in \textbf{Appendix D}. 
We compare G-CARL with the corresponding base models (Base), supervised fine-tuning (SFT), and an MLLM-as-a-Judge reward baseline. 
Additionally, we evaluate a diverse range of general-purpose and medical LVLMs under a zero-shot inference setting.

\begin{table*}[!t]
\centering
\caption{Main results on our medical report interpretation benchmark. Numbers are mean over 3 seeds with standard deviation. }
\vspace{-8pt}
\label{tab:main}
\setlength{\tabcolsep}{9.2pt}
\renewcommand{\arraystretch}{0.92}
\scalebox{0.85}{
\begin{tabular}{l|cccc|cc}
\toprule
\multirow{2}{*}{\textbf{Method}}
& \multicolumn{4}{c|}{\textbf{Subjective Metrics}}
& \multicolumn{2}{c}{\textbf{Objective Metrics}} \\
\cmidrule(lr){2-5}\cmidrule(lr){6-7}
& \textbf{Overall} & Accuracy & Satisfaction & Expression & Precision (\%) & Recall (\%) \\
\midrule
% =========== Closed-source reference ===========
\rowcolor{rowgray}
\multicolumn{7}{l}{\textit{General LVLMs}} \\
GPT-4o~\cite{hurst2024gpt}          
& 1.588$_{\pm0.007}$ & 0.974$_{\pm0.007}$ & 0.358$_{\pm0.001}$ & 0.256$_{\pm0.001}$ & 95.26 & 42.63 \\
ERNIE 4.5 VL~\cite{ernie45technicalreport}          
& 1.709$_{\pm0.010}$ & 1.088$_{\pm0.012}$ & 0.376$_{\pm0.002}$ & 0.245$_{\pm0.000}$ & 95.27 & 53.74 \\
% 1.740 -》 1。709
GLM-4.6V~\cite{glm46v}              
& 1.819$_{\pm0.014}$ & 1.159$_{\pm0.013}$ & 0.395$_{\pm0.002}$ & 0.265$_{\pm0.001}$ & 95.90 & 59.51 \\
Step-3.7-Flash~\cite{step37flashblog} 
& 1.842$_{\pm0.011}$ & 1.175$_{\pm0.012}$ & 0.412$_{\pm0.002}$ & 0.255$_{\pm0.001}$ & 95.96 & 72.11 \\
Kimi K2.5~\cite{kimi2.5}            
& 1.903$_{\pm0.002}$ & 1.221$_{\pm0.002}$ & 0.418$_{\pm0.001}$ & 0.264$_{\pm0.001}$ & 97.63 & 76.42 \\
% \rowcolor{rowpink}
\rowcolor{rowours}
Gemini 3.1 Pro~\cite{gemini31pro}
& \textbf{1.964$_{\pm0.008}$} & \textbf{1.229$_{\pm0.008}$} & \textbf{0.438$_{\pm0.003}$} & \textbf{0.297$_{\pm0.002}$} & \textbf{98.10} & \textbf{80.29} \\
\midrule
\rowcolor{rowgray}
\multicolumn{7}{l}{\textit{Medical LVLMs}} \\
Hulumed~\cite{hulumed}         
& 1.020$_{\pm0.018}$ & 0.480$_{\pm0.013}$ & 0.282$_{\pm0.004}$ & 0.258$_{\pm0.002}$ & 72.28 & 32.65 \\
Medgemma~\cite{sellergren2026medgemma} 
& 1.004$_{\pm0.005}$ & 0.460$_{\pm0.008}$ & 0.289$_{\pm0.004}$ & 0.255$_{\pm0.001}$ & 74.30 & 33.86 \\
Lingshu~\cite{lingshu}
& \textbf{1.527$_{\pm0.003}$} & \textbf{0.902$_{\pm0.002}$} & \textbf{0.357$_{\pm0.003}$} & \textbf{0.268$_{\pm0.002}$} & \textbf{89.69} & \textbf{39.41} \\
\midrule
\rowcolor{rowgray}
\multicolumn{7}{l}{\textit{Qwen-VL series}} \\
% % =========== Qwen-VL series ===========
% Qwen3-VL-2B-Instruct (Base)
% & 1.323$_{\pm0.019}$ & 0.732$_{\pm0.020}$ & 0.344$_{\pm0.003}$ & 0.246$_{\pm0.001}$ & -- & -- \\
% \quad +SFT
% & 1.324$_{\pm0.014}$ & 0.704$_{\pm0.016}$ & 0.365$_{\pm0.003}$ & 0.255$_{\pm0.001}$ & -- & -- \\
% \quad +MLLM-as-a-Judge          & --$_{\pm--}$ & --$_{\pm--}$ & --$_{\pm--}$ & --$_{\pm--}$ & --$_{\pm--}$ & --$_{\pm--}$ \\
% \rowcolor{rowours}
% \quad \textbf{+Ours}         & --$_{\pm--}$ & --$_{\pm--}$ & --$_{\pm--}$ & --$_{\pm--}$ & --$_{\pm--}$ & --$_{\pm--}$ \\
% --- Qwen 4B ---
Qwen3-VL-4B-Instruct (Base)
& 1.527$_{\pm0.010}$ & 0.897$_{\pm0.007}$ & 0.374$_{\pm0.003}$ & 0.256$_{\pm0.001}$ & 92.19 & 49.35 \\
\quad +SFT
& 1.603$_{\pm0.009}$ & 0.945$_{\pm0.012}$ & 0.392$_{\pm0.002}$ & 0.266$_{\pm0.002}$ & 92.90 & 57.28 \\
\quad +MLLM-as-a-Judge
& 1.662$_{\pm0.012}$ & 0.993$_{\pm0.009}$ & 0.396$_{\pm0.002}$ & 0.273$_{\pm0.001}$ & 93.82 & 57.72 \\
\rowcolor{rowpink}
\quad \textbf{+Ours}
& \textbf{1.709$_{\pm0.009}$} & \textbf{1.028$_{\pm0.008}$} & \textbf{0.406$_{\pm0.000}$} & \textbf{0.275$_{\pm0.002}$} & \textbf{93.97} & \textbf{58.92} \\
% 0.267--》0.275
% --- Qwen 8B ---
Qwen3-VL-8B-Instruct (Base)
& 1.626$_{\pm0.012}$ & 0.980$_{\pm0.010}$ & 0.388$_{\pm0.006}$ & 0.258$_{\pm0.001}$ & 93.46 & 60.68 \\
\quad +SFT
& 1.718$_{\pm0.008}$ & 1.040$_{\pm0.010}$ & 0.403$_{\pm0.002}$ & 0.275$_{\pm0.001}$ & 94.22 & 63.44 \\
\quad +MLLM-as-a-Judge
& 1.766$_{\pm0.008}$ & 1.089$_{\pm0.007}$ & 0.407$_{\pm0.002}$ & 0.270$_{\pm0.001}$ & 95.85 & 65.47 \\
\rowcolor{rowpink}
\quad \textbf{+Ours}
& \textbf{1.829$_{\pm0.009}$} & \textbf{1.141$_{\pm0.008}$} & \textbf{0.411$_{\pm0.001}$} & \textbf{0.277$_{\pm0.001}$} & \textbf{96.62} & \textbf{72.18} \\
% 0.265--》0.277
% =========== InternVL2.5 series ===========
\midrule
\rowcolor{rowgray}
\multicolumn{7}{l}{\textit{InternVL3 series}} \\
InternVL3-8B (Base)
& 1.358$_{\pm0.007}$ & 0.744$_{\pm0.007}$ & 0.345$_{\pm0.001}$ & 0.269$_{\pm0.001}$ & 91.25 & 40.37 \\
\quad +SFT
& 1.513$_{\pm0.018}$ & 0.868$_{\pm0.016}$ & 0.385$_{\pm0.004}$ & 0.260$_{\pm0.001}$ & 92.06 & 50.88 \\
\quad +MLLM-as-a-Judge   
& 1.553$_{\pm0.018}$ & 0.897$_{\pm0.017}$ & 0.388$_{\pm0.001}$ & 0.271$_{\pm0.001}$ & 92.39 & 52.61 \\
\rowcolor{rowpink}
\quad \textbf{+Ours}         
& \textbf{1.638$_{\pm0.008}$} & \textbf{0.967$_{\pm0.010}$} & \textbf{0.397$_{\pm0.001}$} & \textbf{0.274$_{\pm0.001}$} & \textbf{92.51} & \textbf{57.97} \\
% --- InternVL 26B ---
InternVL3-14B (Base)
& 1.616$_{\pm0.016}$ & 0.986$_{\pm0.013}$ & 0.365$_{\pm0.004}$ & 0.265$_{\pm0.001}$ & 92.93 & 44.86 \\
\quad +SFT
& 1.663$_{\pm0.007}$ & 0.995$_{\pm0.008}$ & 0.400$_{\pm0.000}$ & 0.268$_{\pm0.001}$ & 93.34 & 58.58 \\
\quad +MLLM-as-a-Judge          
& 1.670$_{\pm0.003}$ & 1.010$_{\pm0.004}$ & 0.399$_{\pm0.003}$ & 0.261$_{\pm0.000}$ & 93.72 & 58.06 \\
\rowcolor{rowpink}
\quad \textbf{+Ours}         
& \textbf{1.739$_{\pm0.007}$}
& \textbf{1.061$_{\pm0.005}$}
& \textbf{0.405$_{\pm0.000}$}
& \textbf{0.273$_{\pm0.002}$}
& \textbf{95.42}
& \textbf{64.57} \\
\bottomrule
\end{tabular}}
\end{table*}

% We compare G-CARL against zero-shot inference (Base), supervised fine-tuning (+SFT), and standard GRPO (+MLLM-as-a-Judge) across multiple model scales. 

% \begin{table}[t]
% \centering
% \setlength{\tabcolsep}{4.8pt}
% \renewcommand{\arraystretch}{0.95}
% \caption{External medical QA evaluation on public Chinese medical benchmarks. CMB and Chinese-SimpleVQA are used only as external sanity checks and are not involved in training, reward construction, model selection, or hyperparameter tuning.}
% \label{tab:external_medical_qa}
% \scalebox{0.86}{
% \begin{tabular}{lccc}
% \toprule
% Method & CMB-Train Acc. & CMB-Val Acc. & Chinese-SimpleVQA-Test Acc. \\
% \midrule
% Base MLLM & -- & -- & -- \\
% SFT & -- & -- & -- \\
% MLLM-as-a-Judge RL & -- & -- & -- \\
% \rowcolor{ourrow}
% \textbf{G-CARL} & \textbf{--} & \textbf{--} & \textbf{--} \\
% \bottomrule
% \end{tabular}}
% \vspace{-6pt}
% \end{table}

\begin{table}[t]
\centering
\setlength{\tabcolsep}{6 pt}
\renewcommand{\arraystretch}{0.9}
\caption{Evaluation results on the CMB dataset.}
\vspace{-8pt}
\label{tab:cmb_evaluation}
\scalebox{0.83}{
\begin{tabular}{lccc}
\toprule[1.5pt]
\multirow{2}{*}{Method}
& \multicolumn{2}{c}{QA Accuracy (\%)}
& \multicolumn{1}{c}{Open-ended Generation} \\
\cmidrule(lr){2-3}\cmidrule(lr){4-4}
& Train & Val & Professionalism \\
\midrule
Base            & 74.85 & 71.42 & 3.58 \\
SFT             & 75.00 & 71.43 & 3.48 \\
MLLM-as-a-Judge & 74.96 & 70.36 & 3.51 \\
\textbf{G-CARL} & \textbf{75.48} & \textbf{72.05} & \textbf{3.61} \\
\bottomrule[1.5pt]
\end{tabular}}
\vspace{-6pt}
\end{table}

% CMB and Chinese-SimpleVQA are used only as external sanity checks and are not involved in training, reward construction, model selection, or hyperparameter tuning.CMB and Chinese-SimpleVQA are used only as external sanity checks and are not involved in training, reward construction, model selection, or hyperparameter tuning.


\begin{table}[t]
\centering
% \setlength{\tabcolsep}{5.5pt}
\renewcommand{\arraystretch}{0.85}
\caption{Ablation study of reward designs in G-CARL.}
\vspace{-8pt}
\label{tab:tab2}
\setlength{\tabcolsep}{2.8pt}
\scalebox{0.83}{
\begin{tabular}{lcccccc}
\toprule[1.5pt]
    & \textbf{Overall} & Acc & Sat & Exp & Pre & Rec \\
\midrule
\multicolumn{7}{l}{Reward Combination} \\
\midrule
Baseline
    & 1.626 & 0.980 & 0.388 & 0.258 & 93.46 & 60.68 \\
+ $R_\text{check}$ + $R_\text{format}$
    & 1.720 & 1.046 & 0.406 & 0.268 & 95.36 & 64.10 \\
+ $R_\text{fact}$ + $R_\text{format}$
    & 1.748 & 1.074 &0.404 & 0.270 & 95.84 & 62.81 \\
+ $R_\text{fact}$ + $R_\text{check}$
    & 1.810 & 1.135 & 0.411 & 0.264 & 96.03 & 70.09 \\
+ $R_\text{fact}$ + $R_\text{check}$ + $R_\text{format}$
    & \textbf{1.829} & \textbf{1.141} & \textbf{0.411}
    & \textbf{0.277} & \textbf{96.62} & \textbf{72.18} \\
\midrule
\multicolumn{7}{l}{Reward Design}     \\
\midrule
$R_\text{check}$: w/ static rubric
    & 1.786 & 1.100 & 0.412 & 0.274 & 94.83 & 66.09 \\
$R_\text{fact}$: w/o retrieval
    & 1.734 & 1.057 & 0.406 & 0.271 & 94.26 & 63.81      \\
\bottomrule[1.5pt]
\end{tabular}}
\vspace{-18pt}
\end{table}


\subsection{Main Results on MMedReport}
\noindent \textbf{Quantitative Comparison.}
Table~\ref{tab:main} compares G-CARL with general-purpose LVLMs, specialized medical LVLMs, and different training paradigms based on the Qwen3-VL and InternVL3 backbones. 
While the MLLM-as-a-Judge reward improves over SFT, its holistic rubric struggles to jointly optimize medical accuracy, demand satisfaction, and expression quality. 
By decomposing reward supervision into externally verifiable medical claims and internally grounded checklist objectives, G-CARL achieves the highest scores across both objective and subjective metrics. 
On Qwen3-VL-8B, G-CARL improves the overall subjective score, while boosting claim-level precision (+0.77\%) and checklist-level recall (+6.71\%), 
indicating more informative and richer interpretations. 
Moreover, G-CARL consistently outperforms specialized medical LVLMs and remains competitive with substantially larger general-purpose LVLMs under zero-shot evaluation, demonstrating the effectiveness of G-CARL.

\begin{figure*}[t]
    \centering
\includegraphics[width=0.95\linewidth]{exp1.pdf}
\vspace{-8pt}
    \caption{Qualitative comparison of interpretation results across different models. Two representative real-world cases (Case A and Case B) are presented, where Case B includes both report images and dialogue history as input.}
    \vspace{-10pt}
    \label{fig:visualize}
\end{figure*}
% Our model demonstrates superior medical accuracy, demand satisfaction, and expression quality.

\vspace{1pt}
\noindent \textbf{External Generalization Study.}
We further transfer the trained models to CMB without any adaptation. As shown in Table~\ref{tab:cmb_evaluation}, G-CARL improves QA accuracy on both splits (+0.63) and attains the highest professionalism score in open-ended generation (3.61), whereas SFT and MLLM-as-a-Judge bring marginal gains and even degrade professionalism. This indicates that grounding rewards in verifiable medical evidence suppresses hallucinated content and better elicits the medical accuracy already latent in the base model, rather than merely fitting the PMRI response style.


\vspace{1pt}
\noindent \textbf{Case Study.}
% 找几个报告单的例子进行分析
% We further examine the complementary effects of the two reward components. 
% The claim-verification reward reduces hallucinations by penalizing unsupported medical claims, but optimizing it alone often yields overly conservative responses with insufficient coverage of patient concerns. 
% Conversely, the dynamic checklist reward improves demand satisfaction by encouraging coverage of key findings and user-specific needs, yet alone it tends to increase verbosity and may reward unnecessary elaboration. 
% As shown in Figure~\ref{fig:reward_reliability}, the two rewards induce opposite biases in response length and content coverage. 
% Their joint optimization balances evidence-bounded factuality and patient-centered completeness, leading to responses that are both faithful and sufficiently informative.
As shown in Fig. \ref{fig:visualize}, we qualitatively compare the outputs of our method with SFT and MLLM-as-a-Judge using Qwen3VL-8B as the base model. The report shows a normal pH (7.432), elevated chloride (111.0 mmol/L), and severe anemia (Hb = 58 g/L). However, SFT incorrectly diagnoses metabolic acidosis and misclassifies the severe anemia as mild, while MLLM-as-a-Judge mistakes the elevated chloride level for hypochloremia instead of hyperchloremia. In contrast, our method correctly identifies all key abnormalities and produces clinically accurate interpretations. Similarly, in the right example, SFT fails to incorporate the user's smoking and alcohol cessation history, while MLLM-as-a-Judge fails to address the user's primary concern of whether to continue the medication. 
\begin{figure}[h]
    \centering    \includegraphics[width=0.99\linewidth]{human_preference_final1.pdf}
    \vspace{-8pt}
    \caption{Human preference evaluation by clinicians and patients. "Ours/B", "Ours/S", and "Ours/M" denote pairwise comparisons between G-CARL and the base model, the SFT model, and the MLLM-as-a-Judge-trained model.}
    \vspace{-18pt}
    \label{fig:preference}
\end{figure}
% Clinicians noted that G-CARL responses are less prone to hallucinated advice and more directly address the patient's core anxieties compared to the generic elaborations produced by standard RL.
In contrast, our method effectively integrates historical context, directly answers the user's question, and provides evidence-grounded recommendations.


% \noindent \textbf{Human Preference Evaluation.}
% To further validate the reliability of our method, we conduct a blind pairwise human preference study on 250 held-out real-world cases. Three professional clinicians independently evaluate anonymized response pairs generated by G-CARL and the baselines across the three evaluation dimensions as well as overall preference, with final judgments determined by majority voting. In addition, we recruit 50 participants without medical training to assess the comprehensibility of the generated interpretations.
% As shown in Fig.~\ref{fig:preference}, G-CARL is consistently preferred over both SFT and MLLM-as-a-Judge. Notably, our method achieves clear gains in medical accuracy and demand satisfaction, with preference margins of 122:72 and 106:64, respectively, over the MLLM-as-a-Judge baseline. It also receives substantially higher comprehensibility ratings from non-expert participants, indicating that the generated interpretations are easier for patients to understand. We emphasize that the patient study complements rather than replaces expert clinical assessment.

\noindent \textbf{Human Preference Evaluation.}
We conduct a blind pairwise human preference study on 250 held-out cases. Three professional clinicians compare anonymized responses generated by G-CARL and the baselines across the three evaluation dimensions and overall preference, with majority voting used for the final decision. We additionally recruit 50 participants without medical training to evaluate the comprehensibility of the generated interpretations.
As shown in Fig.~\ref{fig:preference}, G-CARL is consistently preferred over both SFT and MLLM-as-a-Judge, particularly in medical accuracy and demand satisfaction, with preference margins of 136:85 and 106:64, respectively. It also achieves substantially higher comprehensibility ratings from non-expert participants, suggesting that its responses are easier for patients to understand.



% \begin{table}[t]
% \centering
% \setlength{\tabcolsep}{5.9pt}
% \renewcommand{\arraystretch}{0.95}
% \caption{Ablation study on the reward components.}
% \label{tab:ablation_study}
% \scalebox{0.82}{
% \begin{tabular}{l|cccccc}
% \toprule
% Method & Acc & Rel & Exp & Overall & Precision & Recall \\
% \midrule
% Baseline     & 0.980 & 0.388 & 0.258 & 1.626 & 93.46 & 60.68 \\
% w/o R_{check}                         & 1.046 & 0.406 & 0.268 & 1.720 & -- & -- \\
% w/o R_{fact}                       & 1.100 & 0.412 & 0.274 & 1.786 & -- & -- \\
% w/o R_{format}                         &   &  &  &  & -- & -- \\
% \rowcolor{ourrow}\textbf{Ours} & \textbf{1.141} & \textbf{0.411} & \textbf{0.277} & \textbf{1.829} & -- & -- \\
% \bottomrule
% \end{tabular}}
% \vspace{-10pt}
% \end{table}


\subsection{Ablation Study}
\noindent \textbf{Effect of Reward Components.} 
% To evaluate the contribution of each reward component, we conduct ablation studies on Qwen3VL-8B. As shown in Table~\ref{tab:tab2}, introducing $R_\text{check}$ together with $R_\text{format}$ consistently improves the baseline, indicating that checklist-based supervision primarily enhances coverage of case-specific requirements. In contrast, incorporating $R_\text{fact}$ substantially improves medical factuality. Combining all three rewards achieves the best performance across all evaluation metrics, demonstrating that they provide complementary supervision. We further analyze the key reward designs. Replacing the case-specific checklist with a static rubric consistently reduces overall performance and checklist coverage, confirming the importance of adapting supervision to individual consultation contexts. Likewise, removing retrieval from the factuality reward noticeably degrades overall performance and claim-level precision, demonstrating that retrieval-grounded claim verification provides more reliable factual supervision than relying solely on the model's parametric knowledge.
To evaluate the contribution of each reward component, we conduct ablation studies on Qwen3VL-8B. 
As shown in Table~\ref{tab:tab2}, $R_\text{check}$ combined with $R_\text{format}$ improves the baseline by enhancing case-specific requirement coverage, while $R_\text{fact}$ provides substantial gains in medical factuality. 
Combining all three rewards achieves the best performance, demonstrating their complementary effects. 
Further analysis shows that replacing the dynamic checklist with a static rubric degrades overall performance, highlighting the importance of case-adaptive supervision. 
Similarly, removing retrieval from $R_\text{fact}$ reduces claim-level precision, confirming the effectiveness of retrieval-grounded factual verification.

% We further examine the two key reward designs. Replacing the case-specific checklist with a static rubric consistently degrades performance, confirming the importance of adapting supervision to individual consultation contexts. Likewise, removing retrieval from the factuality reward leads to a noticeable performance drop, showing that retrieval-grounded claim verification provides more reliable factual supervision than relying solely on the model's parametric knowledge.



% In contrast, RGCV yields larger gains, improving Acc from 1.040 to 1.100 and Overall from 1.718 to 1.786, demonstrating the importance of claim-level verification. Combining both components achieves the best performance, showing that RGCV and RDCS are complementary for improving overall interpretation quality.


% \begin{figure}[t]
%     \centering
% \includegraphics[width=0.9\linewidth]{AnonymousSubmission/LaTeX/grpo_training_curves.png}
%     \caption{Human preference evaluation by professional clinicians. G-CARL shows significant win rates against baselines in overall response quality.}
%     \label{fig:preference}
% \end{figure}

% \vspace{1pt}
% \noindent \textbf{Comparison with Other RL Methods.}
% We further compare G-CARL with representative RL methods for open-ended generation in Table~\ref{tab:tab3}. G-CARL consistently achieves the best overall performance across all evaluation dimensions. Preference-based methods (e.g., DPO) rely heavily on high-quality preference data, limiting their effectiveness for the diverse response space of PMRI. Factuality-oriented methods (e.g., FactScore and CAPRL) primarily optimize factual correctness, while rubric-based methods (e.g., PROMETHEUS and RAR) lack explicit evidence verification. In contrast, G-CARL combines retrieval-grounded claim verification with objective-specific reward decomposition, jointly optimizing medical factuality, demand satisfaction, and expression quality. Compared with MedRepBench, which  emphasizes structured clinical finding recall, G-CARL produces more accurate, patient-centered, and clinically reliable interpretations.

\vspace{1pt}
\noindent \textbf{Comparison with Other RL Methods.}
Table~\ref{tab:tab3} compares G-CARL with representative RL methods for open-ended generation. Preference-based DPO performs worst, since high-quality preference data can hardly cover the diverse response space of PMRI. PROMETHEUS scores responses directly against reference answers, which is too coarse-grained to yield confident judgments and thus provides unstable reward signals. Rubric-based RAR offers finer criteria but lacks explicit evidence verification, while MedRepBench emphasizes structured clinical finding recall and is therefore less aligned with user-specific demands. Factuality-oriented rewards (FactScore, CapRL) bring the most pronounced accuracy gains among the baselines. By coupling retrieval-grounded claim verification with objective-specific reward decomposition, G-CARL attains the best overall performance, jointly improving three dimensions.


% \begin{table}[t]
% \centering
% \setlength{\tabcolsep}{2.2pt}
% \renewcommand{\arraystretch}{0.95}
% \caption{Comparison with existing RL methods.}
% \vspace{-10pt}
% \label{tab:tab3}
% \scalebox{0.85}{
% \begin{tabular}{l|cccc}
% \toprule
% Method & Acc & Rel & Exp & \textbf{Overall} \\
% \midrule
% DPO~\cite{rafailov2023direct} & 0.438 & 0.332 & 0.262 & 1.032 \\
% \midrule
% \multicolumn{5}{l}{GRPO-based methods} \\
% \midrule
% \quad + FactScore~\cite{min2023factscore} & 1.100 & 0.412 & 0.268 & 1.781 \\
% \quad + PROMETHEUS~\cite{kim2024prometheus} & 0.530 & 0.347 & 0.255 & 1.132 \\
% \quad + CapRL~\cite{xing2025caprl} & 1.050 & 0.410 & 0.269 & 1.729 \\
% \quad + RAR~\cite{gunjal2025rubrics}  & 1.016 & 0.398 & 0.269 & 1.682 \\  
% \quad + MedRepBench~\cite{shang2025medrepbench} &1.048 &0.406 &0.268 &1.722 \\
% \quad + \textbf{Ours} & \textbf{1.141} & \textbf{0.411} & \textbf{0.277} & \textbf{1.829} \\
% \bottomrule
% \end{tabular}}
% \vspace{-14pt}
% \end{table}
\begin{table}[t]
\centering
\setlength{\tabcolsep}{2.5pt}
\renewcommand{\arraystretch}{0.92}
\caption{Comparison with existing RL methods.}
\vspace{-10pt}
\label{tab:tab3}
\scalebox{0.83}{
\begin{tabular}{l|cccc}
\toprule[1.5pt]
Method & \textbf{Overall} & Acc & Sat & Exp \\
\midrule
DPO~\cite{rafailov2023direct} & 1.032 & 0.438 & 0.332 & 0.262 \\
\midrule
\multicolumn{5}{l}{GRPO-based methods} \\
\midrule
\quad + PROMETHEUS~\cite{kim2024prometheus} & 1.132 & 0.530 & 0.347 & 0.255 \\
\quad + RAR~\cite{gunjal2025rubrics} & 1.683 & 1.016 & 0.398 & 0.269 \\
\quad + MedRepBench~\cite{shang2025medrepbench} & 1.722 & 1.048 & 0.406 & 0.268 \\
\quad + CapRL~\cite{xing2025caprl} & 1.729 & 1.050 & 0.410 & 0.269 \\
\quad + FactScore~\cite{min2023factscore} & 1.776 & 1.100 & 0.408 & 0.268 \\
\quad + \textbf{Ours} & \textbf{1.829} & \textbf{1.141} & \textbf{0.411} & \textbf{0.277} \\
\bottomrule[1.5pt]
\end{tabular}}
\vspace{-14pt}
\end{table}


\section{Conclusion}
In this paper, we present PMRI as a challenging yet underexplored task that requires both evidence-grounded medical factuality and context-dependent patient communication. We propose G-CARL, a reinforcement learning framework that decomposes reward supervision through retrieval-grounded claim verification and case-specific checklist guidance. Extensive experiments on the newly constructed \textsc{MMedReport} benchmark, together with clinician-authored evaluation and human preference studies, show that G-CARL consistently improves the quality of patient-oriented medical report interpretation. We hope this work lays the foundation for more reliable and patient-centered multimodal medical assistants.

% % We present patient-oriented medical report interpretation (PMRI) as a crucial yet underexplored task that demands evidence-bounded factuality and context-dependent patient communication. Identifying the limitations of existing reward formulations, we proposed G-CARL, a verifiability-aware reinforcement learning framework. By decoupling the reward signal into retrieval-grounded claim verification and reference-guided dynamic checklists, G-CARL provides fine-grained and robust supervision. 
% Extensive experiments on the newly constructed \textsc{MMedReport} benchmark, supported by a clinician-authored evaluation protocol and human preference studies, demonstrate that G-CARL significantly improves response quality.
% To address the heterogeneous verifiability of these objectives, we propose \textbf{G-CARL}, a reinforcement learning framework that decomposes reward supervision into retrieval-grounded claim verification and case-specific checklist guidance. By aligning each objective with an appropriate supervision strategy, G-CARL provides more effective and interpretable reward signals for reinforcement learning. Extensive experiments on the newly constructed \textsc{MMedReport} benchmark, together with clinician-authored evaluation and human preference studies, demonstrate that G-CARL consistently improves the quality of patient-oriented medical report interpretation. We hope this work provides a practical step toward more reliable and patient-centered multimodal medical assistants.

\bibliography{aaai2027}

% Check whether the conference requires a reproducibility checklist to be included in the paper.
% If so, you can uncomment the following line and ajust the path to include it.
% \input{ReproducibilityChecklist.tex}

\end{document}
